\documentclass[10pt,a4paper]{amsart}
\usepackage[margin=19mm]{geometry}
\usepackage{amsmath,amssymb,mathtools}
\usepackage[scr=rsfs,cal=euler]{mathalfa}
\usepackage{xfrac}
\usepackage[unicode,hidelinks,bookmarks=false]{hyperref}
\newcommand{\ie}{i.e.\ }
\newcommand{\cf}{cf.\ }
\newcommand{\resp}{respectively\ }
\newcommand{\Subsection}{\subsection}
\providecommand{\nobreakdash}{\nobreak\hbox{-}}
\setlength{\emergencystretch}{2em}
\multlinegap0pt
\usepackage{bm}
\usepackage{bbm}
\usepackage{dsfont}
\usepackage{xspace}
\usepackage{cancel}
\usepackage{mathtools}
\usepackage{amsthm}
\makeatletter
\@ifundefined{theorem}{\newtheorem{theorem}{Theorem}}{}
\@ifundefined{proposition}{\newtheorem{proposition}[theorem]{Proposition}}{}
\makeatother
\title[Clustering, order conditions, and languages of interval exch]{Clustering, order conditions, and languages of interval exchanges}
\author{S\'ebastien Ferenczi, Luca Q. Zamboni}
\date{Source: arXiv:2507.17370v2 (CC BY 4.0)}
\begin{document}
\maketitle

\noindent\emph{Source and licence.} Clustering, order conditions, and languages of interval exchanges. Version arXiv:2507.17370v2 (CC BY 4.0), distributed under the Creative Commons Attribution 4.0 International licence, as recorded in the arXiv licence metadata for this exact version and shown on the abstract page https://arxiv.org/abs/2507.17370v2. Full rights notice: licenses/arxiv-2507.17370-CC-BY-4.0.txt.\par\smallskip

\noindent\textit{Editorial scope.} Counted unit: the Proposition whose source label is \texttt{pfl2} in the article (source0.tex lines 641--643, source label \texttt{pfl2}), delivered together with the article's own complete original proof of it (source0.tex lines 644--692, one \texttt{proof} environment of 49 lines). No mathematics is written, repaired or re-derived: statement and proof are the article's own text line for line. Required context retained uncounted: pfl (lines 577--579). Every definition, display and label of those lines is retained unchanged. Omitted from this package and not redistributed: 1--576, 580--640, 693--877. Nothing inside a retained excerpt is omitted or altered: every retained body is delivered exactly as the article's lines are written, so no omission receipt of any kind accompanies this package. Citations: the retained excerpts carry no citation command at all, so no bibliography entry of the article is transcribed and no reference is redistributed. Reference adaptation: none was needed and none was made; every reference inside the delivered unit resolves inside it, so no unresolved reference or citation is produced. Printed slips kept literal: no printed word, symbol, hypothesis or number of the counted statement or of its proof was altered. Layout and class: the article's own class is \texttt{amsart} with packages including color, times, amsfonts, amsthm, amsmath, psfrag, latexsym, amssymb, graphicx, epsfig, euscript, tikz, xy, caption; the wrapper is the lane's pinned \texttt{amsart} editorial wrapper at 10pt on A4, which restates the theorem-like environments the retained text needs and adds the article's own macro definitions that the retained text needs (none), with extra wrapper packages (bm, bbm, dsfont, xspace, cancel, mathtools). The delivered numbering is the unit's own. Assets: no retained span contains a figure environment, a graphics-inclusion command, a PDF-page inclusion, a table or a TikZ picture; no image file is copied into this package, the package declares no asset dependency and no image file is redistributed with it. Licence: version arXiv:2507.17370v2 of this article is distributed under the Creative Commons Attribution 4.0 International licence, as recorded in the arXiv licence metadata for this exact version and shown on the abstract page https://arxiv.org/abs/2507.17370v2. A full rights notice is delivered at licenses/arxiv-2507.17370-CC-BY-4.0.txt. Characters: every retained excerpt is pure ASCII, so the delivered engine, XeLaTeX with its default Latin Modern text font, reports no missing character.
\begin{proposition}\label{pfl}
Let $W$ fe a finite set of words. The non extendable language $F(W)$ satisfies the order condition for  $(<_D,<_A)$ if and only if $W$ is included in an (extendable) language satisfying the order condition for $(<_D,<_A)$, with no connection.
\end{proposition}

\begin{proposition}\label{pfl2}
The non extendable language $F(W)$ satisfies the order condition for  $(<_D,<_A)$ if and only if $W$ is included in a grouped coding of
an affine interval exchange transformation with the same pair of  orders.\end{proposition}

\begin{proof}
The ``if" direction is immediate. 

In the other direction, let $\mathcal A$, of cardinality $k$, be the set of all letters in  $F(W).$ Starting from $F(W)$, we build the $W_n$ as in Proposition \ref{pfl}, but only   when $n\leq N$, $N$ being  the maximal length of the words in $W$.  We have  $\mathcal A=\{b_1<_D\ldots <_D b_k\}$,  $\mathcal A=\{a_1<_A\ldots <_A a_k\}$. \\
 Note that if $w$ is in $W_{n-2}$  and we look at its extensions in $W_n$, it can be followed by $b_j,,\ldots ,b_{j+m}$ for some $j$ and $m,$, preceded by $a_i, \ldots ,a_{i+h}$ for some $i$ and $h$ (if there is a gap, this contradicts the order condition and saturation at a  previous stage). \\


We recall  some properties of words of length $n$ in the language $L(T)$ for an interval exchange transformation $T$ with no connection. Such a word $w$ corresponds to  the cylinder $[w]$ for $T$, which is an interval $J_w=I_{w_1}\cup T^{-1}I_{w_2}.\cup \ldots \cup T^{-n+1}I_{w_n}$ whose endpoints are among the $T^p\gamma_i$, $0\leq p\leq n-1$.   
A word $w$ of  length $n$ is left 
special whenever $J_w$ contains a $\beta_l$, right special whenever $J_w$ contains a $T^{-n}\gamma_l$. Suppose  $w$ is bispecial; it can be preceded by $a_i<_A\ldots <_A a_{i+h}$, and followed by $b_j<_D\ldots <_D b_{j+m}$. Then $J_w$ contains $T^{-n}\gamma_l$, $j \leq l\leq j+m-1$
and $\beta_l$, $i \leq l\leq i+h-1$. By the geometry of $T$, $a_ib_j$ is in $L(T)$; then $a_ib_{j+1}$ is in $L(T)$ if $T^{-n}\gamma_j<\beta_i$, $a_{i+1}b_j$ is in $L(T)$ if $T^{-n}\gamma_j>\beta_i$. Continuing from left to right, we get from the resolution of $w$ a unique (strict) ordering
 of the concerned $T^{-n}\gamma_l$ and $\beta_l$ (for the usual order, left to right, on the interval); we call it the $(L(T), w)$-ordering. It can be described by integer intervals $Z_l$, $i-1\leq l\leq i+h$, some of them possibly empty, with $l_1<l_2$ when $l_1\in Z_l$, $l_2\in Z_{l+1}$,
 whose disjoint union is the integer interval $[j,j+m-1]$, such  that $\beta_l<T^{-n}\gamma_{l'}< \beta_{l+1}$ when $l'$ is in $Z_l$, $i \leq l\leq i+h-1$, $T^{-n}\gamma_{l'}< \beta_{i}$ when $l'$ is in $Z_{i-1}$, $T^{-n}\gamma_{l'}> \beta_{i+h}$ when $l'$ is in $Z_{i+h}$.
 
 
 Let $w$ be a bispecial word in $W_{n-2}$, and we look at its extensions  in $W_n$; then  it can be followed by $b_j,\ldots ,b_{j+m}$ for some $j,m$, preceded by $a_i, \ldots ,a_{i+h}$ for some $i,h$ (if there is a gap, this contradicts the order condition and saturation at a  previous stage). If $L$ was an $L(T)$ for an interval exchange transformation $T$ with no connection, $W_n$ would define an $(L(T),w)$ ordering described by  integer intervals $Z_l$, $i-1\leq l\leq i+h$, as above.
 This ordering depends only on $W_n$ because of the saturation of the extension graph (there has to be an edge form $a_i$ to $b_j$, then either an edge from $a_i$ to $b_{j+1}$ or an edge from $a_{i+1}$ to
$b_j$, 
and so  on), we call it the $(W_n,w)$ ordering.  \\


 We shall now build interval exchange transformation $T_n$ so that the set of words of  length $n$ of  $L(T_n)$ is $W_n$. We consider them  as generalised interval exchange transformations with the same defining intervals: indeed, the $T_n$  are piecewise affine on their (common) defining intervals, and are also affine interval exchange transformation but with a growing number  of defining intervals.  We shall denote by $J_{n,w}$ the interval corresponding to the cylinder $[w]$ for $T_n$. At the beginning we choose
points $0=\gamma_{0}<.\ldots < \gamma_{k}=1$, define  $I_{b_l}=[\gamma_l,\gamma_{l+1})$ for $0\leq  l \leq k-1$.

Then we  choose  points $0=\beta_0<\ldots < \beta_{k}=1$, such that the natural ordering (from left to right) of the points  $\gamma_l$, $1\leq l\leq k-1$
and $\beta_l$, $1\leq l\leq k-1$ is the  same as the $(W_2, \varepsilon)$ ordering. We  define $T_2$ as the affine  interval exchange transformation for which  $I_{b_l}=[\gamma_l,\gamma_{l+1})$ for $0\leq i\leq k-1$, and $T_2I_{a_l}=[\beta_l,\beta_{l+1})$ for $0\leq i\leq k-1$. Then the set of words of length $2$ of $L(T_2)$ is $W_2$. 

To build $T_3$, we look at words of length $1$ and their bilateral extensions in $W_3$. If $w$ is not bispecial, we put $T_3=T_2$ on $J_{2,w}$. If $w$ is bispecial, preceded by $a_i<_A\ldots <_A a_{i+h}$, and followed by $b_j<_D\ldots <_D b_{j+m}$, its extensions in $L(T_2)$ are determined by the $(L(T_2),w)$ ordering
 while its extensions in $W_3$ are determined by the $(W_3, w)$ ordering, given by integer intervals $Z_l$ as above, and these may give different extensions. So we shall modify $T_2$ into $T_3$ such that  the $(W_3, w)$ ordering and the $(L(T_3), w)$ ordering  give the same extensions; for this, we choose arbitrarily some new points
  $\gamma_{i,3}$ in $J_{2,w}$, 
 such that 
 \begin{itemize}
 \item $\gamma_{l,3}<\gamma_{l',3}$ whenever $T_2^{-1}\gamma_l<T_2^{-1}\gamma_{l'}$,
 \item   $\beta_l<\gamma_{l',3}< \beta_{l+1}$ when $l'$ is in $Z_l$, $i\leq l\leq i+h-1$, $\gamma_{l',3}< \beta_i$ when $l'$ is in $Z_{i-1}$, $\gamma_{l',3}> \beta_{i+h}$ when $l'$ is in $Z_{i+h}$.
\end{itemize}
 When $l$ is such that $T_2^{-1}\gamma_l$ is in no $J_{2,w}$ for any bispecial $w$, we put $\gamma_{l,3}=T_2^{-1}\gamma_l$. Thus we  have added $k-1$ points $\gamma_{i,3}$. And on $J_{2,w}$ we replace the affine $T_2$ by  a piecewise affine $T_3$  such that  $T_3J_{2,w}=T_2J_{2,w}$ 
  and $T_3\gamma_{i,3}=\gamma_l$; thus  the set of words of length $3$ of $L(T_3)$ is $W_3$. \\

Similarly, suppose we have built $T_{n-1}$. To build $T_n$, we look at words of length $n-2$, in $W_{n-2}$  or equivalently in  $L(T_{n-2})$ or  $L(T_{n-1})$,  and their bilateral extensions in $W_n$. If $w$ is not bispecial, we put $T_n=T_{n-1}$ on $J_{n-1,w}$. If $w$ is bispecial, preceded by $a_i<_A\ldots <_A <a_{i+h}$, and followed by $b_j<_D\ldots <_D b_{j+m}$, we modify $T_{n-1}$ into $T_n$ such that  the $(W_n, w)$ ordering, given by the integer intervals $Z_l$, and the $(L(T_n), w)$ ordering  give the same extensions; for this, we create new points $\gamma_{i,n}$ in $J_{n-1,w}$, such 
that 
 \begin{itemize}
 \item $\gamma_{l,n}<\gamma_{l',n}$ whenever $T_{n-1}^{-1}\gamma_{l,n-1}<T_{n-1}^{-1}\gamma_{l',n-1}$ 
  \item   $\beta_l<\gamma_{l',n}< \beta_{l+1}$ when $l'$ is in $Z_l$, $i\leq l\leq i+h-1$, $\gamma_{l',n}< \beta_i$ when $l'$ is in $Z_{i-1}$, $\gamma_{l',n}> \beta_{i+h}$ when $l'$ is in $Z_{i+h}$.
\end{itemize}
When $l$ is such that $T_{n-1}^{-1}\gamma_{l,n-1}$ is in no $J_{n-1,w}$ for any bispecial $w$, we  put $\gamma_{l,n}=T_{n-1}^{-1}\gamma_{l,n-1}$. And on $J_{n-1,w}$ we replace $T_{n-1}$, which is affine on this set, by  a piecewise affine $T_n$  such that 
 $T_nJ_{n-1,w}=T_{n-1}J_{n-1,w}$  and $T_n\gamma_{l,n}=\gamma_{l,n-1}$. Thus we  have added $k-1$ points $\gamma_{l,n}$, and, by the induction hypothesis  that $T_{n-1}^{n-3}\gamma_{l,n-1}=\gamma_l$, we have   $T_n^{-n+2}\gamma_l=\gamma_{l,n}$, which ensures that  the set of words of length $n$ of $L(T_n)$ is $W_n$. 


Thus the proposition is proved with the interval exchange transformation $T_N$. \end{proof}

\end{document}
